 \documentclass[11 pt, a4paper]{article}
 %%Language and font encodings
\usepackage[utf8]{inputenc}
\usepackage[english, bulgarian]{babel}
\usepackage[T1, T2A]{fontenc}
%%Sets page size and margins
\usepackage[a4paper, top=3cm, bottom=3cm,left=3cm,right=3cm,marginparwidth=1.75pt]{geometry}
%%Useful packages
\usepackage{caption}
\usepackage{textcomp}
\usepackage{verbatim}
\usepackage{color, listings}
\usepackage[subrefformat=parens]{subcaption}
\usepackage[toc,page]{appendix}
\usepackage[usenames,dvipsnames,svgnames]{xcolor}
\usepackage{algorithmicx}
\usepackage{algorithm}
\usepackage{amsfonts}
\usepackage{amssymb}
\usepackage{amsthm}
\usepackage{gensymb}
\usepackage{mathrsfs}
\usepackage{array}
\usepackage{bm}
\usepackage{booktabs}
\usepackage{colortbl}
\usepackage{enumitem}
\usepackage{fancyvrb}
\usepackage{placeins}
\usepackage{standalone}
\usepackage{subcaption}
\usepackage{tabulary}
\usepackage{textcomp}
\usepackage{tikz}
\usepackage{url}
\usepackage{varwidth}
\usepackage{indentfirst}
\usepackage{wrapfig}
\usepackage{amsmath}
\usepackage{systeme}
\usepackage{graphicx}
\usepackage[colorinlistoftodos]{todonotes}
\usepackage[colorlinks=true, allcolors=black]{hyperref}
\usepackage{enumitem}
\usepackage{fmtcount}
\usepackage{multicol}
\usepackage{breqn}
\setlength{\columnsep}{1cm}
\setlength{\parindent}{2em}
\renewcommand{\baselinestretch}{1.75}

% Lemma, Corollary and Proof to show off.
\newtheorem{thm}{Теорема}
\newtheorem{lm}{Лема}
\newtheorem{cl}{Следствие}

\theoremstyle{definition}
\newtheorem{defin}{Дефиниция}
\newtheorem{p}{Задача}
\newtheorem*{e}{Пример}

\theoremstyle{remark}
\newtheorem*{s}{Решение}
\begin{document}
\begin{center}
\pagenumbering{gobble}
\vspace{- 1.5cm}
\large\textbf {ДЕВЕТНАДЕСЕТА УЧЕНИЧЕСКА КОНФЕРЕНЦИЯ НА \\ УЧЕНИЧЕСКИЯ ИНСТИТУТ ПО МАТЕМАТИКА И ИНФОРМАТИКА\\УС'19}\\
\vspace{3cm}
\huge\textbf {ПРИЛОЖЕНИЯ НА ВЕРОЯТНОСТНИЯ МЕТОД}\\
\vspace{3cm}
\large{Aвтор:}\\
\Large\scshape Анна Михалкова\\
(СМГ „Паисий Хилендарски“, град София)\\
\vspace{3cm}
\large{Научен ръководител:\\}
\large{\textsc{Димитър Чакъров}\\
(МГ „Акад.К.Попов“, град Пловдив)}
\end{center}
\newpage
\begin{normalsize}
\begin{center}
\textbf{Резюме}
\end{center}
Разработеният реферат е в областта на комбинаториката и съставлява едно широко обобщение на вероятностния метод. Въведени са основната терминология и теоритичния апарат, нужен за прилагане на метода. Към тях са предоставени примерни задачи с директно приложение. Главният фокус е около състезателни задачи от национални и международни
 форуми, в чиито решения метода може да се използва и допринася за олекотяване на решенията. Една от основните ни цели е показване широката приложимост на вероятностния метод не само в теоритични клонове на математиката, но и в задачи със състезателен характер. Разгледаните задачи са с повишена трудност и смятаме, че биха били от полза на ученици при подготовката им за олимпиади и състезания по математика.
\end{normalsize}
\begin{normalsize}
\begin{center}
\textbf{Summary}
\end{center}
The following project is in the field of combinatorics. It is a broad summary of the probabilistic method. Basic definitions and theorems, needed for the proper use of the probabilistic method, are introduced. Simple examples are provided in order to illustrate the theory in practice and to introduce the method itself. The project focuses on competitive problems from major national and  international competitions, in whose solutions could be used probabilistic method, contributing to their briefness. One of our main aims is demonstrating the broad application of the probabilistic method not only in some theoretical fields of mathematics, but also in competitive problems. The problems in the abstract are with increased difficulty and we believe, they can serve as a preparation for mathematical olympiads and competititons.

\end{normalsize}
\pagenumbering{arabic}
\newpage
\section{Въведение}
През последните години вероятностният метод се развива интензивно и става едно от най-използваните и широкоразпространени средства за решаване на комбинаторни проблеми. Една от главните причини за голямото развитие в областта е съществената роля на теорията на вероятности в компютърните науки - област, източник на множество интригуващи комбинаторни задачи.\par
Началото на метода поставя Пол Ердьош (един от най-плодовитите математици на 20. век) в края на четиридесетте години на миналия век . Не е изненада и че той има значителен принос за развитието на областта и вида, в който я познаваме.\par
Основната идея на вероятностния метод е да докаже съществуването на математическа структура с определени свойства, от които ние се интересуваме. Разглеждаме произволно състояние на конструкцията с дадени параметри и задаваме вероятност на елементарните събития, изграждащи я. Използвайки това доказваме, че желаните свойства съществуват с положителна вероятност, тоест има състояние, в което търсеното от нас е изпълнено. Истинността на получените чрез метода резултати се предполага от формалната дефиниция за вероятност. Ето защо и той намира широко приложение в теория на числата, алгебра, анализ, теория на информация, комбинаторика, а и косвено в области като алгоритмичното програмиране.\par
В следващата секция на този проект са въведени важни факти, дефиниции и използваната терминология. В третата секция са дадени задачи, за които е нужно единственно директното приложение на основните понятия. Четвъртата секция представя някои от необходимите теореми за прилагане на вероятностния метод. Петата секция предлага още задачи, които са от значително по-голяма трудност и изискват немалка съобразителност по отношение на мястото на метода в цялостното решение. Последните две секции са съответно заключение, бъдещо развитие и благодарности. 
\newpage
\section{Дефиниции и означения}
В тази секция излагаме важни дефиниции и нотацията, която ще използваме в реферата.
\begin{defin}{\textit{Случайна величина}}\\
Случайна величина е такава променлива, която мени стойността си произволно. (Например, ако се разгледат възможните точки при хвърляне на зарче със $6$ страни ($D_6$), събитието да бъдат получени определен брой точки е случайна величина.)
\end{defin}
\begin{defin}{\textit{Вероятност}}\\
Вероятността да се случи някакво елементарно събитие (или още отношението на броя на благоприятните изходи към общия им брой) ще бележим с $P$.\par
Пример:\\
В случая със зарчето\\
$P(D_6=1)=P(D_6=2)=P(D_6=3)=P(D_6=4)=P(D_6=5)=P(D_6=6)=\dfrac{1}{6}$;\\
$P(D_6=0)=0$;\\
$P(D_6\geq4)=\dfrac{1}{2}.$
\end{defin}
\begin{defin}{\textit{Математическо очакване}}\\
Математическо очакване представлява характеристична стойност на вероятностното разпределение на случайна величина. Математическо очакване за случайна величина $X$ може да се интерпретира като нейна средна стойност, въпреки че тази стойност може да не бъде възможен неин изход. Математическото очакване не трябва да се бърка с най-вероятния изход!\par
Формалната дефиниция е
\begin{equation*}
E[X]=\sum_{x}P(X=x) \cdot x 
\end{equation*}
Но нека се върнем на примера със зарчето:
\begin{equation*}
E[D_6]=\dfrac{1}{6}\cdot1+\dfrac{1}{6}\cdot2+ \dfrac{1}{6}\cdot3+\dfrac{1}{6}\cdot4+\dfrac{1}{6}\cdot5+
\dfrac{1}{6}\cdot6=3.5
\end{equation*}
Накратко математическото очакване е сумата от възможните изходи, умножени съответно с вероятността да се случат.
\end{defin}
\begin{thm}[Линейност на математическото очакване]
При дадени случайни величини $X_1,X_2,\dots,X_n$ винаги важи
\begin{equation*}
E[X_1+X_2+\dots+X_n]=E[X_1]+E[X_2]+\dots+E[X_n]
\end{equation*}
Теоремата е очевидна, ако приемем, че събитията са независими едно от друго-ако хвърлим зарче $100$ пъти очакваме да получим средно $350$ точки. По силното твърдение е, че теоремата важи и за събития, които зависят едно от друго. В много от случаите $E[X]$ може да се пресметне при разбиване на главното събитие на елементарни такива и прилагане на теоремата за линейност.\par
\end{thm}
\section{Задачи за загрявка }
В тази секция даваме лесноусвоими примери за директно приложение на вероятностния метод.
\begin{p}[AHSME 1989]

$7$ момичета и $13$ момчета са се наредили в редица. Нека означим с $S$ колко пъти момче и момиче в редицата стоят едно до друго.(Например в редицата $12211121211121211211$, където на мястото на $1$ има момиче, а на $2$-момче, $S=12$). Намерете математическото очакване за стойността на $S$.\par
\end{p}
\begin{s}
Трябва да получим колко общо измежду деветнадесетте двойки са с момче и момиче -- това е същото,като да разгледаме всяка двойка поотделно и да съберем стойностите. Ще означаваме двойка момче-момче или двойка момиче-момиче с $0$, а двойка момче-момиче или момиче-момче с $1$. С $S$ означаваме общия брой на двойките момче-момиче и момиче-момче, а $S_i$($i=1,2,\dots,19$) приема стойност $0$ или $1$ в съвпадаща с тази на вида на двойката. Знаем, че:
\begin{equation*}
S=S_1+S_2+\dots+S_{19}
\end{equation*}
Всичко това ни навежда на мисълта да приложим теоремата за линейност на математическото очакване. Преди това обаче трябва да изчислим вероятността една двойка да е момче-момиче или момиче-момче:\\
\begin{equation*}
P=\dfrac{7\cdot13}{\dfrac{20\cdot19}{2}}=\dfrac{7\cdot13}{10\cdot19}=\dfrac{91}{190}\\
\end{equation*}
\begin{equation*}
E[S]=E[S_1]+E[S_2]+\dots+E[S_{19}]=19\cdot\dfrac{91}{190}=9.1\\
\end{equation*}
\end{s}
\begin{p}[NIMO 5.6]

Том има научен калкулатор. За нещастие, всички копчета са счупени освен $1$,$2$,$3$,$+$ и $-$. Той въвежда поредица от пет знака, като всеки от тях има еднаква вероятност да бъде избран. След това калкулаторът пресмята израза и извежда резултат $Е$.
Намерете математическото очакване за стойността на $Е$.\par
\end{p}
\begin{s}
Ако заменим плюс с минус и обратното, средната стойност на всички числа, с които сме извършили действието, остава числото,образувано от цифрите преди поява на знак. Затова можем да означим плюсовете и минусите със сигнал СПРИ. За всеки израз се зачита стойността му до сигнал СПРИ. Ако разгледаме един израз и заменим $1$ с $3$ и обратно, средната стойност на първоначалния израз и новополучения се записва само с двойки на мястото на цифрите $1$, $2$ и $3$. Следователно можем да заменим всяка цифра с $2$. Извършваме определената размяна.\\
Вероятността един знак да е $2$ е $\dfrac{3}{5}$. Означаваме $p=\dfrac{3}{5}$. Вероятността първият знак да бъде $2$ е $p$, а за всеки следващ става $10^k p$ за $k=1,2,3,4$. Прилагаме теоремата за линейността и получаваме израза:
\begin{equation*}
2(p+10p^2+10^2 p^3+10^3 p^4+10^4 p^5)=2\cdot p\cdot\dfrac{(10p)^5-1}{10p-1}=1866
\end{equation*}

\end{s}
\begin{p}[AIME 2006/6]

Нека $S$ е множеството от реалните числа, които могат да бъдат представени като десетична дроб $0.a b c$, където $a$, $b$ и $c$ са различни цифри. Намерете сумата на елементите на $S$.\par
\end{p}
\begin{s}
В числото $abc$ $a$,$b$,$c$ за различни, следователно броят на елементите на $S$ е $10\cdot9\cdot8=720$. Вероятността едно число да принадлежи на $S$ e $\dfrac{1}{2}$(всяко от числата или удовлетворява условието, или не). Чрез теоремата за линейността, сумата на дробите става $720\cdot\dfrac{1}{2}=360$.\\
\end{s}
\begin{p}[NIMO 4.3]

Един ден кон и офицер са на квадратчета на един и същи ред от безкрайна шахматна дъска. Когато изведнъж започва голяма метеоритна буря, пада метеор във всяко квадратче независимо и случайно с вероятност $p$. Нито конят, нито офицерът са оцелени, но тяхното движение може да бъде възпрепятствано. За каква стойност на $p$ очакваният брой на квадратчетата, до които може да стигне коня за един ход, е равен на очаквания брой квадратчета, до които може да стигне офицерът за един ход?\par
\end{p}
\begin{s}
Всеки кон може да стигне до $8$ полета с един ход, а вероятността едно поле да е без метеорит е $1-p$. Един офицер може да стигне до $4$ полета на разстояние $n$ по диагонал от неговото поле($n$ варира от $1$ до безкрайност). За да може да стигне до поле на разстояние $n$, трябва и всички полета преди избраното, и самото то да не съдържа метеорит. Следователно вероятността да може да стигне до поле на разстояние $n$ от сегашното е $(1-p)^n$.\\
За да намерим стойността, искана в условието, изравняваме двете стойности за математическото очакване за достъпните полета, използвайки теоремата за линейността. Така достигаме до уравнението:\\
\begin{equation*}
    8(1-p)=4((1-p)+(1-p)^2+(1-p)^3+\dots)
\end{equation*}
\begin{equation*}
2(1-p)=((1-p)+(1-p)^2+(1-p)^3+\dots)
\end{equation*}
$p=1$ е решение както на уравнението, така и на задачата. Ако $p$ е различно от $1$:
\begin{equation*}
1=((1-p)+(1-p)^2+\dots)
\end{equation*}
\begin{equation*}
1=\dfrac{(1-p)(1-(1-p)^n)}{p}
\end{equation*}
При нарастване на $n$ до безкрайност $(1-p)^n$ клони към $0$, а съответно $1-(1-p)^n$ към $1$. Оттук
\begin{equation*}
\\lim_{n\to\infty}((1-p)+(1-p)^2+(1-p)^3+\dots)=\dfrac{1-p}{p}
\end{equation*}
\begin{equation*}
1=\dfrac{1-p}{p}
\end{equation*}

Следователно окончателно $p$ е равно на $1$ или $\dfrac{1}{2}$.
\end{s}
%\begin{p}[NIMO 7.3]
%Иван има четири безкрайно големи купчини от монети от $1$, $5$, $10$ и $25$ стотинки. Той избира една монета от някоя купчина и я взема. Процесът се повтаря докато събере цяло число левове. Какво е математическото очакване за него, изразено в стотинки?\par
%\end{p}
%%Нека $x_n$ е състоянието, в което Иван има $n$ стотинки (mod$100$). Ако $n$ е различно от $0$, има вероятност $\dfrac{1}{4}$ да бъде добавена монета от $1$, $5$, $10$ или $25$ стотинки. Знаем, че $x_0=0$. Прилагаме теоремата за линейността и получаваме уравнението:\\ 
 %\begin{equation*}
 %x_0=\dfrac{1}{4}((x_{n+1}+1)+(x_{n+5}+5)+(x_{n+10}+10)+(x_{n+25}+25)) \end{equation*}
%\end{s}
\section{Полезна теория}
В тази секция предоставяме на читателя няколко теореми, които могат да бъдат използвани в по-сложни многостъпкови задачи от състезания и олимпиади.
\begin{lm}
За положителни $n$ и $k$
\begin{equation*}
{{n}\choose{k}}={\dfrac{1}{e}}\bigg(\dfrac{en}{k}\bigg)^k,
\end{equation*}
където $e\approx2.718\dots$е Неперовото число.
\end{lm}
\begin{thm}[Неравенство на Джордж Бул]
Нека $A_1, A_2,\dots,A_k$ са елементарни събития. Ако
\begin{equation*}
P(A_1)+P(A_2)+\dots+P(A_k)<1,
\end{equation*}
то има положителна вероятност никое от събитията да не се случи.
\end{thm}
\begin{thm}[Неравенство на Марков]
Нека $X$ е произволна променлива, която приема само неотрицателни стойности, а $E[X]=c$.  Следователно
\begin{equation*}
P(X\geq rc)\leq \dfrac{1}{r}.
\end{equation*}
\end{thm}
\begin{thm}[Лема на László Lovász]
Съществуват няколко елементарни събития, чиято вероятност е не по-голяма от $p$, такива, че всяко от тях е зависимо от най-много $d$ от останалите събития. Тогава ако
\begin{equation*}
epd\leq1,
\end{equation*}
 вероятността никое от събитията да не се случи е положителна.
\end{thm}
\section{Състезателни задачи}
\begin{p}[Russia 1996]

В парламента има $1600$ народни представители, които формират $16000$ комитети с по $80$ човека всеки. Докажете, че съществуват два комитета, които имат поне $4$ общи члена.\par
\end{p}
\begin{s}
Означаваме с $X$ броя на народните представители, които са членове на два комитета (произволни), а с $X_1, X_2,\dots,X_{1600}$ дали всеки представител поотделно е член и на двата комитета. $X_i$ приема стойност $1$, ако човекът е член, и стойност  $0$, ако не е. Забелязваме, че
\begin{equation*}
X=X_1+X_2+\dots+X_{1600}
\end{equation*}
От това следва, че можем да приложим теорема за линейността:
\begin{equation*}
E[X]=\sum_{i=1}^{i=1600} X_i5
\end{equation*}
С $n_i$ означаваме броя на комитетите, в които i-тият представител членува. От условието намираме, че $\sum_{i=1}^{i=1600}n_i=80.16000=1280000$, а средната стойност на $n_i=\dfrac{16000.80}{1600}=800$. Също така знаем, че
\begin{equation*}
E[X_i]=\dfrac{{{n_i}\choose{2}}}{{{16000}\choose{2}}}.
\end{equation*}
От горните факти можем да направим следната оценка
\begin{equation*}
E[X]\geq1600\cdot{800\choose{2}}/{16000\choose{2}}=1600.\dfrac{800\cdot799}{16000\cdot15999}=3.995
\end{equation*}
и тъй като ${E[X]}$ по условие трябва да бъде цяло, доказателството е извършено.
\end{s}
\begin{p}
Покажете, че може да се организира кръгов турнир (всеки играе с всеки по веднъж) с повече от $1000$ човека, така че във всяка група от $1000$ човека може да се намери такъв, който е победил всички от групата си.\par
\end{p}
\begin{s}
Нека $n$ е общият брой на участниците, а една група се състои от $k$ участника (в случая $k=1000$). Тогава ${n}\choose{k}$ е броят на всички такива групи. Вероятността да няма участник, който побеждава всички останали e $(1-\dfrac{1}{2^k})$. Прилагайки Теорема 4, за да съществува турнир с исканите свойства, трябва да е изпълнено:
\begin{equation*}
{{n}\choose{k}}\Big(1-\dfrac{1}{2^k}\Big)^{n-k}<1
\end{equation*}
Пресмятаме за $k=1000$ и използвайки лемата, достигаме до:
\begin{equation*}
{{n}\choose{1000}}\Big({\dfrac{2^{1000}-1}{2^{1000}}}\Big)^{n-1000}<1
\end{equation*}
\begin{equation*}
{{n}\choose{1000}}\Big({\dfrac{2^{1000}-1}{2^{1000}}}\Big)^{n-1000}<\dfrac{1}{e}\Big(\dfrac{en}{1000}\Big)^{1000}\Big({\dfrac{2^{1000}-1}{2^{1000}}}\Big)<1,
\end{equation*}
което е вярно за достатъчно голямо $n$.
\end{s}
\begin{p}[BAMO 2004]

Дадени са $n$ реални числа, не всички равни на нула, но  със сума равна на нула. Докажете, че числата могат да бъдат означени $a_1, a_2,\dots a_n$, така че\\
\begin{equation*}
a_1 a_2+a_2 a_3+\dots+a_n a_1<0.
\end{equation*}
\end{p}
\begin{s}
Да допуснем, че всички суми са по-големи или равни на нула. Произведението $a_i a_j$ ($i$ е различно от $j$) се появява $n(n-2)!$ пъти (тук стойностите на $a_i$ и $a_j$ не се променят-фиксирали сме ги).\\
\begin{equation*}
    n(n-2)!\sum a_i a_j\geq0
\end{equation*}
Следователно
\begin{equation*}
2\sum a_i a_j\geq0
\end{equation*}
и 
\begin{equation*}
\sum_{i=1}^{i=n} a_i^2+2\sum a_i a_j>0
\end{equation*}
От друга страна горепосоченият израз е точно равен на сумата на елементите на квадрат, което е $0$. Така достигаме до противоречие с допуснатото и твърдението е доказано.\par
\end{s}
\begin{p}[Russia 1999]

В училище всяко момче харесва поне едно момиче. Докажете, че съществува подмножество $S$ от поне половината ученици, така че всяко момче харесва нечетен брой момичета от $S$.\par
\end{p}
\begin{s}
Ще покажем как такова множество може да бъде построено. Първо ще разпределим момичетата -- за всяко от тях има вероятност да бъде избрано $\dfrac{1}{2}$ (или ще бъде, или няма да бъде в подмножеството). Така избрани броят на харесваните от определено момче момичета в групата вече може да се намери -- ако е нечетен, включваме момчето в подмножеството, а ако е четен -- не. Следователно вероятността да бъде избрано едно момче също е $\dfrac{1}{2}$. Така от теоремата за линейността получаваме, че полученото множество се състои от поне половината ученици, с което доказателството е завършено.
\end{s}
\begin{p}[IMC 2002]

На олимпиада иа $6$ задачи и $200$ участника, които са много умни и затова всяка задача е решена от поне $120$ от тях. Докажете, че съществуват двама участника, за които всяка задача е решена от поне един от двамата.\par
\end{p}
\begin{s}
Съществуват ${200\choose{2}}=19900$ различни двойки участници.
Всяка задача не е решена от максимум ${80\choose{2}}=3160$ двойки участници. Следователно вероятността да не бъде решена в определена двойка е $\dfrac{3160}{19900}$. Изчисляваме математическото очакване за броя на двойките, в които има нерешена задача. Означаваме със $Z$ броя на нерешените задачи, а $Z_1, Z_2,\dots,Z_6$ са променливи, които показват дали всяка задача е решена или не от поне един от двойката (приемат стойност $1$ при нерешена задача и $0$ при решена). Знаем, че:
\begin{equation*}
Z=Z_1+Z_2+\dots+Z_6
\end{equation*}
\begin{equation*}
E[Z]=E[Z_1]+\dots+E[Z_6]=6\cdot\dfrac{3160}{19900}=\dfrac{18960}{19900}<1
\end{equation*}
Тъй като математическото очакване за нерешените задачи е нецяло число по-малко от едно, а в действителност броят на задачите е цяло число, то със сигурност съществува двойка, която разполага с $0$ нерешени задачи, тоест всяка задача е решена от поне един от двамата.\\
\end{s}
\begin{p}[Shortlist 1999 C4]

Нека $А$ е множество от $N$ остатъка по модул $N^2$. Покажете, че съществува множество $B$ от $N$ остатъка по модул $N^2$, така че 
$A+B=\{a+b|a\in A,b\in B\}$ съдържа поне половината остатъци по модул $N^2$.\par
\end{p}
\begin{s}
Означаваме с $X$ броя на всички възможни остатъци, които $A+B$ съдържа. Нека остатъкът $i\in \{A+B\}$. Тъй като $|A|=n$, съществуват $n$ на брой остатъка $b\in B$, които удовлетворяват $A+b\equiv i(\mod{N^2})$. Вероятността остатъкът $i$ да принадлежи на $A+B$ e $1-\bigg(1-\dfrac{n}{n^2}\bigg)^n$От всички възможни изходи изваждаме неблагоприятните изходи-всеки от възможните $n$ остатъци не принадлежи на множеството $B$. Така определяме и математическото очакване за броя на остатъците:
\begin{equation*}
E[X]=n^2\cdot\bigg(1-\bigg(1-\dfrac{n}{n^2}\bigg)^n\bigg)
\end{equation*}
За да се съдържат поне половината остатъци, трябва единствено да докажем:
\begin{equation*}
1-\bigg(1-\dfrac{n}{n^2}\bigg)^n\geq\dfrac{1}{2}
\end{equation*}
\begin{equation*}
\dfrac{1}{2}\geq \bigg(1-\dfrac{n}{n^2}\bigg)^n.
\end{equation*}
Последното неравенство е тривиално и може да бъде доказано по индукция. \\
\end{s}
\begin{p}[Iran TST 2008/6]

Нека $799$ отбора играят в турнир всеки срещу всеки по веднъж. Докажете, че съществуват две групи $A$ и $B$, всяка с по $7$ отбора, които нямат общ отбор, така че всички отбори от едната група са победили всички от другата.\par
\end{p}
\begin{s}
Означаваме с $А$ произволна група от $7$ отбора. Нека $X$ е броят на отборите, които са победени и от седемте отбора в група $А$. Построяваме насочен граф с върхове отборите и ребра изиграните мачове. Ребрата са насочени от победител към победен. С $d_v$ означаваме ребрата, които завършват във връх $v$. От условието знаем, че $\sum_v d_v={{799}\choose{2}}$, от което заключаваме, че средната стойност на $d_v$ е $399$. Изчисляваме $E[X]$:
\begin{equation*}
E[X]=\sum_v{d_v\choose{7}}/{799\choose{7}}\geq 799\cdot{399\choose{7}}/{799\choose{7}}\approx 800\cdot\bigg(\dfrac{1}{2}\bigg)^7\approx 6,25
\end{equation*}
Тъй като математическото очакване за победените отбори е нецяло число, а в действителност те са цяло число, оценката е достатъчна и със сигурност съществуват $7$ победени отбора от отборите в група $А$ (ще ги наречем група $B$).
\end{s}
\begin{p}[Romania 2004]

Докажете, че за които и да било комплексни числа $z_1, z_2,\dots,z_n$, удовлетворяващи условието
\begin{equation*}
|z_1|^2+|z_2|^2+\dots+|z_n|^2=1,\end{equation*} могат да се изберат $\epsilon_1,\epsilon_2,\dots,\epsilon_n\in \{-1,1\}$, така че\\
\begin{equation*}
\left|\sum_{k=1}^{k=n} \epsilon_k z_k\right|\leq1.
\end{equation*}\par
\end{p}
\begin{s}
Избираме $\epsilon_i$ произволно, а лявата страна на исканото означаваме с LHS. За всяко $z$ важи $|z|^2=z\overline{z}\bar{}$. Знаем, че:
\begin{equation*}
LHS^2=\sum_{k}\left|z_k\right|^2+2\sum_{i<j}\epsilon_i\epsilon_j(z_i\overline{z_j}\bar{}+z_j\overline{z_i}\bar{})
\end{equation*}
Избираме $\epsilon_i$ и $\epsilon_j$ произволно и равнопоставено, значи в два от случаите произведението им е $-1$, а в другите два е $1$, което ни води до $E(\epsilon_i \epsilon_j)=0$. Следователно след като съберем почленно всички леви страни, сумите с множител $2$ пред тях се съкращават (има равен брой с минус и с плюс пред тях). Получаваме:  
\begin{equation*}
E[LHS^2]=\sum_k\left|z_k\right|^2=1.
\end{equation*}
Това означава, че понякога можем да изберем $\epsilon_i$, така че $LHS^2\leq1$, към което се целим.
\end{s}
\begin{p}[USAMO 2012/6]

Нека $x_1,x_2,\dots,x_n$ са реални числа за $n\geq2$ ($n$ е цяло), удовлетворяващи равенствата $x_1+x_2+\dots+x_n=0$ и $x_1^2+x_2^2+\dots+x_n^2=1$. За всяко подмножество $A\subseteq\{1,2,\dots,n\}$ дефинираме
\begin{equation*}
S_A=\sum_{i\in A} x_i.
\end{equation*}
(Ако $A$ е празно множество, $S_A=0$). Докажете, че за положително число $\lambda$, броят на множествата $A$, за които $S_A\geq\lambda $, е най-много $2^{n-3}/\lambda^2$. Кога се достига равенство?\par
\end{p}
\begin{s}Знаем, че
\begin{equation*}
2\sum_{1\leq i\leq j\leq n}x_i x_j=\bigg(\sum_{i=1}^{n}x_i\bigg)^2-\sum_{i=1}^{n}{x_i}^2=-1
\end{equation*}
Следователно
\begin{equation*}
\sum_{1\leq i\leq j\leq n}x_i x_j=-\dfrac{1}{2}
\end{equation*}
Оттук заключаваме, че 
\begin{equation*}
\sum_{A\subseteq N}{S_A^2}=2^{n-1}\sum_{i=1}^{n}{x_i}^2+2^{n-1}\sum_{i=1}^{n}x_i x_j=2^{n-1}-2^{n-2}=2^{n-2}
\end{equation*}
$S_A^2\geq0$. Забелязваме, че $S_A+S_{N|A}=0$ и оттук най-много половината от $\left|S_A\right|\geq\lambda$.\par
От гореспоменатото директно следва(от неравенство на Марков), че максималният брой на търсените множества е $2^{n-3}/\lambda^2$. За да се достигне максималната стойност, трябва да се достига равенство и при $\left|S_A\right|\geq\lambda$. Следователно $S_A\in \big\{-\lambda,0,\lambda \big\}$. Измежду множествата може да съществува най-много едно със стойност $\lambda$ и едно със стойност $-\lambda$(в противен случай ще бъде нарушено условието за сбора). Така за останалите $n-2$  множества остава стойност $0$ и за тях важи следното:
\begin{equation*}
\sum_{i=1}^{n}x_i^2=2\lambda^2=1
\end{equation*}
От условието $\lambda$ е положително число и заради това $\lambda=\dfrac{1}{\sqrt{2}}$.
\end{s}
\begin{p}[Online Math Olimpiad, Ray Li]
Кевин има $2^n-1$ бисквитки, всяка от които се отличава от другите чрез различно непразно подмножество от $\big\{1,\dots,n\big\}$. Всеки ден той избира на случаен принцип една от все още неизядените бисквитки и изяжда нея и всички други, които са с подмножество на нейното множество. Пресметнете математическото очакване за броя дни, през които Марк е ял бисквитки, докато накрая всички свършат.\par
\end{p}
\begin{s}
Броят на дните е равен на броя на случаите, в които една бисквитка е елиминирана или избрана. Нека $C$ е множеството на бисквитките, които са избрани по време на процеса, $A$ -- множеството, характеризиращо определена бисквитка, $D$ -- броят на дните, а $S=\big\{1,\dots,n\big\}$. $X$ определяме като индикатор($0$ или $1$) за това дали $A\in C$. Ясно е, че
\begin{equation*}
E[D]=E[\left|C\right|]=\sum_{A\subseteq S}E[X(A\in C)]=\sum_{A\subseteq S}P(A\in C)
\end{equation*}
За всяко $A\in C$, $А$ е елиминирана в момента, когато надмножество $A'$ на $A$ е избрано. Следователно
\begin{equation*}
P(A\in C)=P(A=A')=\dfrac{1}{K}=\dfrac{1}{2^{n-\left|A\right|}}
\end{equation*}
$K$-брой на надмножествата на $A$\\
Заместваме в първото равенство:
\begin{equation*}
\sum_{A\subseteq S}P(A\in C)=\sum_{A\subseteq S}2^{\left|A\right|-n}=\sum_{1\leq k\leq n} 2^k\cdot2^{k-n}=\dfrac{4^{n+1}-4}{3\cdot2^n}
\end{equation*}
\end{s}
\begin{p}
Нека $n$ e положително цяло число. С $a_k$ означаваме броя на пермутациите на $n$ елемента, от които $k$ са фиксирани. Пресметнете:
\begin{equation*}
a_1+4a_2+9a_3+\dots+n^2a_n
\end{equation*}
\end{p}
\begin{s}
За произволна пермутация $X$ е броят на фиксираните елементи, a търсеният израз е еквивалентен на $n!E[X^2]$.
\begin{equation*}
E[X^2]=1\cdot\dfrac{a_1}{n!}+2^2\cdot\dfrac{a_2}{n!}+\dots+n^2\cdot\dfrac{a_n}{n!}=\dfrac{a_1+4a_2+\dots+n^2a_n}{n!}
\end{equation*}Лесно се вижда, че 
\begin{equation*}
X^2=X+2\cdot{X\choose{2}}
\end{equation*}
\begin{equation*}
E[X^2]=E[X]+E\bigg[2\cdot{X\choose{2}}\bigg]=E[X]+2\cdot E\bigg[{X\choose{2}}\bigg]
\end{equation*}
Ще изчислим $E[X]$ и $E[{X\choose{2}}]$:
\begin{equation*}
E[X]=\dfrac{1}{n!}\cdot n\cdot(n-1)!=1
\end{equation*}
\begin{equation*}
E\bigg[{X\choose{2}}\bigg]=\sum_{ij}P(i,j=const)={n\choose{2}}\cdot\dfrac{1}{n}\cdot\dfrac{1}{n-1}=\dfrac{1}{2}
\end{equation*}
$P(i,j=const)$-вероятността да съществуват две фиксирани числа $i$ и $j$ сред всички $n$
Следователно
\begin{equation*}
E[X]+E[2\cdot{X\choose{2}}]=n!E[X^2]=n!\cdot(1+2\cdot\dfrac{1}{2})=2n!
\end{equation*}
\end{s}
\begin{p}[Ердьош]

Докажете, че във всяко множество $S$ от $n$ различни положителни цели числа съществува подмножество $T$ от $\left[\frac{1}{3}\right]n$ или повече елемента със свойството $a+b \neq c$ за всеки $a, b, c\in T$ (не непременно различни).\par
\end{p}
\begin{s}
Нека $p=3k+2$ е просто, което удовлетворява неравенството $p>2.\underset{1 \leq i \leq n}{\max}|b_i|$, а $C$ е такова множество: $C=\{k+1, k+2,\dots,2k+1\}$. В $C$ няма два елемента със сбор, който принадлежи на множеството. Забелязваме, че $\dfrac{|C|}{p-1}=\dfrac{k+1}{3k+1}>\dfrac{1}{3}$. Избираме произволно число $x$, такова че $1\leq x<p$ и дефинираме $d_1, d_2,\dots,d_n$, така че $d_i \equiv x b_i (\mod{p})$ , $0\leq d_i<p$. За всяко фиксирано $i$($1\leq i\leq n$) $x$ варира измежду $1,2,\dots,p-1$, a $d_i$ варира измежду ненулевите елементи на множеството $S$. Изчисляваме вероятността $d_i$ да принадлежи на $C$: $P(d_i\in C)=\dfrac{|C|}{p-1}>\dfrac{1}{3}$. Следователно математическото очакване за броя на елементите $b_i$, за които $d_i\in C$, е по-голямо от $n/3$. Това означава, че съществува $x$, $1\leq x<p$ и подредица $А$, принадлежаща на $B$ с дължина $|A|>\dfrac{n}{3}$, така че $x a(\mod{p})\in C$, за всяко $a\in C$. В $А$ със сигурност съществуват някакви $a_1,a_2,a_3$: $a_1+a_2=a_3$ и $x a_1+x a_2\equiv x a_3 (\mod{p})$. От последното равенство следва, че в $C$ има два елемента със сбор, който принадлежи на множеството-противоречие с твърдението в началото на решението.   \\
\end{s}
\begin{p}[Sperner]

Даденo е множеството $A$ от $N$ различни подмножества $S_1,S_2,\dots,S_N$ от $\{1,2,\dots,n\}$, така че никое $S_i$ не е подмножество на $S_j$. Докажете, че
\begin{equation*}
N \leq {{n} \choose {\left[\frac{1}{2}n\right]}}
\end{equation*}
\end{p}
\begin{s}
Означаваме $b$ като произволна пермутация, а $X$ дефинираме като $\big\{i:b(1),\dots,b(i)\in A\big\}$.
Разглеждаме математическото очакване за $X$.
От дефиницията на $A$ по условие променливата $X$ е не по-голяма от $1$, а събитията $\big\{b(1),\dots,b(k)\big\}$ са независими за всяко $k$.
Нека $N_k$ е броят на подмножествата с $k$ елемента в $А$.
\begin{equation*}
    E[X]=\sum_{k=1}^{n}P\left[\big\{b(1),\dots,b(k)\in A\big\}\right]=\sum_{k=1}^{n}\dfrac{N_k}{{n\choose{k}}}\geq\sum_{k=1}^{n}\dfrac{N_k}{{n\choose{[n/2]}}}
\end{equation*}
Тъй като $E[X]\leq1$, то $\sum_{k=1}^{k=n} N_k=N\leq{{n\choose{\left[n/2\right]}}}$, което трябваше да докажем. 
\end{s}
%\begin{p}

%Даден е граф G с множество от върхове V. Докажете, че съществува  независимо множество с най-малко
%\begin{equation*}
%\sum_{v\in V} \left[\frac{1}{deg v+1}\right]
%\end{equation*}
%елементи.\par

%\textit{Решение:}\\
%Независимо множество е множество от върхове на граф, никои два от които са съседни и няма ребро между тях. Да приемем, че сме избрали напълно произволна подредба на V. Означаваме с $I$ независимото множество с най-малък брой елементи. Нека $X_v$ e произволна променлива, която определя дали върхът $v$ принадлежи на множеството $I$. Знаем, че:
%\begin{equation*}
%X=\sum_{v\in V}X_v=\left|I\right|
%\end{equation*}
%и за всяко $v$
%\begin{equation*}
 %E[X_v]=P(v\in I)=\dfrac{1}{deg_v+1}   
%\end{equation*}
%Тогава
%\begin{equation*}
%E[X]=\sum_{v\in V}\dfrac{1}{deg_v+1}
%\end{equation*}
%\end{p}
%и можем да заключим, че съществува специфична подредба със
%\begin{equation*}
%\left|I\right|\geq\sum_{v\in V}\dfrac{1}{deg_v+1}   
%\end{equation*}

\begin{p}
$11n$ точки са разположени по окръжност и са оцветени в $n$ различни цвята, като има по $11$ точки от всеки цвят. Докажете, че могат да се изберат $n$ точки по една от всеки цвят, така че никои две не са съседни.\par
\end{p}
\begin{s}
Означаваме точките по окръжността с $a_1,\dots,a_{11n}$. Вземаме по една точка от всеки цвяти конструираме произволно множество $B$ от $n$ точки. За всяко $1\leq i\leq 11n$ дефинираме събитието $X_i$, когато $a_i$ и $a_{i+1}$ принадлежат на $B$. $X=X_1\cup X_2 \cup\dots\cup X_{11n}$ е обединението на всички възможни неблагоприятни изходи, които искаме да избегнем. Да отбележим, че $P(X_i)=\dfrac{1}{11}\cdot\dfrac{1}{11}=\dfrac{1}{121}$, ако $a_i$ и $a_{i+1}$ са разноцветни, а $P(X_i)=0$ при едноцветни точки. Прилагаме Теорема 3: 
\begin{equation*}
P(X)\leq\sum_{i=1}^{11n}P(X_i)<\dfrac{11n}{11^2}=\dfrac{n}{11}
\end{equation*}
Теоремата обаче не ни дава $P(X)<1$ при $n>11$. Забелязваме, че $X_i$ е независимо от $X_j$ за всички $j$, освен ако $a_i$ или $a_{i+1}$ не е в същия цвят като $a_j$ или $a_{j+1}$. $10$ други точки освен $a_i$ са оцветени в същия цвят и всяка от тях е част от две последователни двойки от точки. Следователно има $2\cdot10+1=21$ двойки $(a_j,a_{j+1})$, различни от$(a_i,a_{i+1})$, които имат един и същ цвят с $a_i$. Аналогично има $21$ двойки, които имат общ цвят с $a_{i+1}$. Така $X_i$ и $X_j$ са независими за всички без най-много $42$ стойности на $j$. Време е да използваме Теорема 4 с вероятност $p=\dfrac{1}{11^2}$ и $d=42$ :
\begin{equation*}
ep(d+1)=e\cdot\dfrac{1}{11^2}\cdot43<\dfrac{28}{10}\cdot\dfrac{43}{121}<1
\end{equation*}
Резултатът след прилагането на теоремата е, че вероятността никой от неблагоприятните изходи да не възникне е положителна. Това означава, че със сигурност има множество от $n$ несъседни две по две точки, които са с различни цветове.  
\end{s}

\begin{p}[Алон и Спенсър, глава 1]
Да се докаже, че за всеки две независими случайни реални величини $X$ и $Y$
\begin{equation*}
P(|X-Y|\leq2)\leq 3P(|X-Y|\leq1)
\end{equation*}
\end{p}
\begin{s}
Нека $\big\{x_i\big\}_{i\leq n}$ са $n$ на брой реални числа, а с $\#$ означаваме брой. Тогава първо ще докажем
\begin{equation*}
\#\big\{(i,j):|x_i-x_j|\leq2\big\}\leq3\# \big\{(i,j):|x_i-x_j|\leq 1\big\}
\end{equation*}
Ще докажем твърдението по индукция. Базата е тривиална.  Допускаме, че твърдението е вярно за всяко множество с по-малко или равно на $n$ точки. Прекарваме синьо ребро между две точки, ако те са на разстояние $1$ или по-малко, и с червено ребро, ако те са на разстояние между $2$ включително и $1$. За целите на броенето ще свържем $x_i$ със себе си с половин синьо ребро, тъй като $(i,i)$ се брои като една двойка, а двойката $(i,j)$ се брои и като $(j,i)$. Целта е да покажем, че червените ребра са най-много два пъти повече от сините ребра.\\
Нека съществуват $n+1$ реални числа $\big\{x_i\big\}$ и нека $y_1<y_2<\dots<y_k$ е максимално дългото подмножество, за което всяка двойка $(y_i,y_{i+1})$ е свързана със червено ребро. Ние твърдим, че при премахване на цялата редица със съответстващите и ребра ще сме премахнали най-много два пъти повече червени от сини ребра. Това е вярно поради следното:\\
Съществуват $k-1$ червени ребра между $(y_i,y_{i+1})$ и съответно $k/2$ сини ребра между $(y_i,y_i)$. За всяко червено ребро от $x_i$ до $y_j$ разглеждаме два случая: $x_i<y_j$ и $x_i>y_j$. Ако $x_i<y_j$, тъй като редицата е максимална, $j>1$, трябва $|x_i-y_{j-1}|\leq 1$ и затова съпоставяме червено ребро $(x_i,y_j)$ на синьо ребро $(x_i, y_{j-1})$. Аналогично ако $x_i>y_j$, следователно $j<k$ и червеното ребро $(y_j,x_i)$ съответства на синьо ребро $(x_i,y_{j+1})$. Всяко премахнато синьо ребро може да бъде съпоставено на най-много едно червено ребро, идващо отдясно, и на най-много едно червено ребро, идващо отляво. Така се оказва, че броят на сините ребра, принадлежащи на нашата редица $y_1,y_2,\dots,y_k$, е най-много половината от броят на червените ребра. И така сме готови с нашето доказателство по индукция.\\

Нека $X_1,X_2,\dots,X_n$ са $n$ независими променливи от условието. Съчетавайки горното твърдение и теоремата за линейността, получаваме, че
\begin{equation*}
E[\#\big\{(i,j):|X_i-X_j|\leq 2\big\}]\leq 3E[\#\big\{(i,j):|X_i-X_j|\leq 1\big\}]
\end{equation*}
\begin{equation*}
n+2{{n}\choose 2}P[|X-Y|\leq 2]\leq 3n+6{{n}\choose 2}P[|X-Y|\leq 1]
\end{equation*}
Границата се достига, когато $n$ клони към безкрайност.
\end{s}

\begin{p}[Алон и Спенсър, глава 1]
Нека $\big\{(A_i,B_i), 1\leq i\leq h \big\}$ е семейство от двойки подмножества на множеството на целите числа, така че $|A_i|=k$ за всяко $i$, $|B_i|=l$ за всяко $i$, $A_i\cap B_i= \emptyset$ и $(A_i\cap B_j)\cup(A_j\cap B_i)\neq \emptyset$ за всяко $i\neq j$. Докажете, че $h\leq\dfrac{(k+l)^{k+l}}{k^{k}l^{l}}$.
\end{p}
\begin{s}
Нека $S$ е множеството от целите числа, които принадлежат на някое от $A_i$ и $B_i$. Приемаме, че $f$ е произволна функция, която изпраща $S$ във $[1,k+l]$, a $X_i$ е събитието, което изпраща $A_i$ във $[1,k]$, a $B_i$-в [k+1, k+l]. Следователно
\begin{equation*}
 P[X_i]=\dfrac{k^k l^l}{(k+l)^{k+l}}   
\end{equation*}
Нека $h>\dfrac{(k+l)^{k+l}}{k^k l^l}$. Тогава математическото очакване за $X_i$ е по-голямо от $1$ и така поне две от събитията се случват с положителна вероятност. Ако случаят е за $X_i$ и $X_j$, тогава $A_i\cap B_j$ и $A_j\cap B_i$ са празни, противоречие с условието.
\end{s}


\begin{p}[Алон и Спенсър, глава 1]
Нека $n\geq4$ и $H$ е $n$-хиперграф с най-много $\dfrac{4^{n-1}}{3^n}$ ръба. Докажете, че съществува оцветяване на върховете на $H$ в $4$ цвята, така че във всеки ръб се съдържат и четирите цвята. (В един $n$-хиперграф всяко ребро, с други думи множество от върхове,съдържа точно $n$ върха.) 
\end{p}
\begin{s}
Произволно оцветяваме всеки връх в един от четирите цвята. Вероятността всеки връх да съдържа най-много $3$ цвята е по-малка или равна на $4\cdot\dfrac{3^n}{4^n}=\dfrac{3^n}{4^{n-1}}$. Така математическото очакване за върховете, които съдържат по-малко от $4$ цвята, е по-малко от $1$. От последното следва, че съществува някакво оцветяване без ръбове с по-малко от $4$ цвята, което и трябваше да докажем. 
\end{s}

%\begin{p}[Алон и Спенсър, глава 1]
%Нека $G=(V,E)$ е граф с $n$ върха и минимална степен $\lambda>10$. Докажете, че върховете $V$ могат да се разделят на две непресичащи се множества $A$ и $B$, така че $|A|\leq O\bigg(\dfrac{n \ln\lambda}{\lambda}\bigg)$ и всеки връх от $B$ има най-малко един съсед от $A$ и най-малко един съсед от $B$.
%\end{p}
%\begin{s}
%Определяме $X$ като произволно подмножество на $G$ , избирайки всеки от върховете на графа с вероятност $p$. Можем да добавим всички върхове $C_X$, които не са съседни до никои от тези в $X$, и $D_X$, чиито единствени съседи са в $X\cup C_X$. Така $A=X\cup C_X\cup D_X$ и всеки връх в $B=V\backslash A$ има поне един съсед в $B$. Пресмятаме, че
%\begin{equation*}
%E[|X|]=np
%\end{equation*}
%\begin{equation*}
%E[|C_X|]=n(1-p)^{\lambda}
%\end{equation*}
%За да намерим мощността на $D_X$, за фиксиран връх $v$ трябва да ограничим възможността, че всички от съседите му принадлежат на $X$ и $C_X$. Означаваме с $a_1,a_2,\dots,a_\lambda$ $\lambda$ от съседите на $v$. Така
%\begin{equation*}
%P[v\in D_X]\leq p^{\lambda}+\lambda(1-p)^{\lambda},
%\end{equation*}
%където $p^{\lambda}$ е вероятността всички съседи да принадлежат на $X$, а с $\lambda(1-p)^\lambda$ ограничаваме вероятността поне един от съседите да принадлежи от $C_X$. Следва, че 
%\begin{equation*}
%E[A]<n(p+\lambda(1-p)^{\lambda}\leq n(p+\lambda e^{-p\lambda})
%\end{equation*}
%\end{s}
\begin{p}[Алон и Спенсър, глава 1]
За всяко естествено число $k$ съществува граф турнир $T_k=(V,E)$ с $|V|>k$, така че за всяко множество $U$ с най-много $k$ върха от $T_k$ съществува връх $v$, за който всички насочени ребра $\big\{(v,u):u\in U\big\}$ принадлежат на $E$. Покажете, че всеки такъв граф турнир съдържа най-малко $\Omega(k2^k)$ върха.
\end{p}
\begin{s}
Нека $T_k$ е граф турнир със свойството $S_k$ (свойството на графа турнир в условието) и с $n$ върха. Избираме произволен връх $v_k$, математическото очакване за чиято степен е по-малко от $n/2$. Фиксираме връх със степен по-малка от $n/2$. Нека $T_{k-1}$ е подтурнирът от върхове, които са победили $v_k$. Тогава $|T_{k-1}|<n/2$ и $T_{k-1}$ има свойство $S_{k-1}$ (Във всяко множество от $T_{k-1}$ с $k-1$ върха, заедно с $v_k$, съществува връх $u$, който побеждава всички останали, включително и $v_k$, и следователно принадлежи на множеството $T_{k-1}$). Продължавайки аналогично, достигаме до $T_k\supset T_{k-1}\supset T_{k-2}\supset\dots\supset T_2\supset T_1$, като за всяко $i$ $T_i$ е множество от върхове, които побеждават $v_k, v_{k+1},\dots, v_i, v_{i+1}$, а $v_i$ е избран връх, така че $2|T_i|<|T_{i+1}|$. Оказва се, че $T_1$ е непразното множество от всички върхове, които побеждават $v_k,\dots,v_2$. Нека $l$ е някой връх извън $T_k$. Заедно $v_k,\dots,v_2,l$ образуват $k$-елементно подмножество на $T_k$ и трябва да има $u\in T_1$, който побеждава всички останали. По принцип обаче няма $l$, който да побеждава всички в $T_1$. И тъй като $T_1$ принадлежи на турнир $T_k$ със свойство $k$, $|T_1|\geq {k+1}$.Така
\begin{equation*}
|T_k|\geq 2^{k-1}|T_1|\geq 2^{k-1}(k+1),
\end{equation*}
което и трябваше да докажем.
\end{s}
\begin{p}[Алон и Спенсър, глава 2]
Нека $n\geq2$ и $H=(V,E)$ е $n$-хиперграф с $|E|=4^{n-1}$ ребра. Докажете, че съществува оцветяване на $V$ с $4$ цвята, така че никое ребро не е едноцветно.(В един $n$-хиперграф всяко ребро, с други думи множество от върхове,съдържа точно $n$ върха.) 
\end{p}
\begin{s}
Оцветяваме върховете на $H$ произволно в четирите цвята. За всяко ребро има вероятност $4^{-n+1}$ да бъде едноцветен. Следователно математическото очакване за едноцветните ребра е $1$. Знаем, че има начин да получим повече от едно едноцветно ребро(всичко е в един цвят) и така заключаваме, че трябва да съществува начин, по който се достига до $0$. 
\end{s}
\begin{p}[Алон и Спенсър, глава 2]
Нека $p>n>10m^2$, $p$ е просто, а $0<a_1<a_2<\dots<a_m<p$ са цели числа. Докажете, че съществува цяло число $x$, $0<x<p$, за което $m$-те на брой числа
\begin{equation*}
((xa_i) (mod p)) mod n, (1\leq i\leq m)
\end{equation*}
са две по две различни.
\end{p}
\begin{s}
Избираме за $x$ някое от $1,2,\dots,p-1$ и дефинираме индикатора $X_{ij}$ като $1$, ако $(x a_i mod p) mod n=(x a_j mod p) mod n$. Съществуват най-много $2(p-1)/n$ кратни на $n$, чиято стойност може да приеме $(x a_i mod p)-(x a_j mod p)$ и всяко от тях е с вероятност $1/p-1$. Следователно математическото очакване за броя на двойките от този тип е най-много
\begin{equation*}
\dfrac{2}{n}{{m}\choose{2}}<1,
\end{equation*}
тъй като $10m^2<n$. От последното заключаваме, че за някоя стойност на $x$ няма две равни числа.
\end{s}
\begin{p}[Алон и Спенсър, глава 2]
Нека $F$ е семейство от подмножествва на множеството $N=\big\{1,2,\dots,n\big\}$ и няма $A,B\in F$, за които $A\subset B$. Тогава с $\sigma\in S_n $ означаваме произволна пермутация на елементите на $N$, а с $X$ случайната величина
\begin{equation*}
X=|\big\{i:\big\{\sigma(1), \sigma(2), \dots, \sigma(i)\big\} \in F\big\}|
\end{equation*}
Пресмятайки математическото очакване за $X$, докажете, че $|F|\leq{{n}\choose{\lfloor{n/2}\rfloor}}$.
\end{p}
\begin{s}
От условието, че няма $A,B\in F$, за които $A\subset B$, знаем, че $X\leq1$. От друга страна всеки елемент $S\in F$ има най-малко вероятност $1/{{n}\choose{\lfloor{n/2}\rfloor}}$ да се появи в $X$. От това и линейността на математическото очакване следва, че $|F|\leq{{n}\choose{\lfloor{n/2}\rfloor}}$.
\end{s}

\begin{p}[Алон и Спенсър, глава 2]
Нека $G=(V,E)$ е двуделен граф с $n$ върха и списък $S(v)$ с повече от $\log_{2}{n}$ цвята за всеки връх $v\in V$. Докажете, че съществува подходящо оцветяване на $G$, така че всеки връх $v$ е в цвят от своя списък $S(v)$. 
\end{p}
\begin{s}
Нека $C$ е множеството, образувано от всички цветове, които се срещат във някой от списъците. За всеки цвят имаме вероятност $\dfrac{1}{2}$ да се срещне в някой от дяловете на графа. Ако всеки връх е в един и същи дял с поне един от цветовете си, твърдението е доказано, като използваме някой от тези цветове. Вероятността за определен връх $v$ всеки от цветовете във списъка му да се среща в другия дял е $\dfrac{1}{2^{|S(v)|}}<\dfrac{1}{n}$. Така матемическото очакване за върховете, които са с несъответстващ цвят, е по-малко от $1$ от линейността на математическото очакване, от което следва и нашият резултат. 
\end{s}

%\begin{p}[Алон и Спенсър, глава 4]
%Нека $X$ е случайна величина, приемаща неотрицателни стойности. С $E(X^2)$ означаваме математическото очакване за нейния квадрат, а с $Var(X)$--дисперсията. Докажете, че
%\begin{equation*}
%P(X=0)\leq \dfrac{Var(X)}{E(X^2}
%\end{equation*}
%\textit{Знаем, че формулата за дисперсия е следната:
%\begin{equation*}
%Var(X)=E(X^2)-[E(X)]^2
%\end{equation*}}
%\end{p}
%\begin{s}
%Нека $p_i=P[X=i]$. Прилагайки Неравенството на Коши-Шварц,
%\begin{equation*}
%(1-p_0)(\sum_{n\geq0}p_{n}n^2)=(\sum_{n\geq1} p_n)(\sum_{n\geq1}p_{n}n^2)\geq(\sum_{n\geq1}p_{n}n)^2=(\sum_{n\geq0}p_{n}n)^2
%\end{equation*}
%Следователно
%\begin{equation*}
%(1-p_0)E[X^2]\geq E[X]^2
%\end{equation*}
%\begin{equation*}
%E[X^2]-E[X]^2\geq p_0E[X^2]
%\end{equation*}
%\begin{equation*}
%Var(X)\geq p_0E[X^2]
%\end{equation*}
%и разделяйки получаваме исканото неравенство.
%\end{s}
%\begin{p}[Алон и Спенсър, глава 4]
%Нека $X$ е случайна величина с математическо очакване $E[X]=0$ и дисперсия $Var(X)=\sigma^2$. Да се докаже, че за всяко $\lambda>0$,
%\begin{equation*}
%P(X\geq \lambda)\leq \dfrac{\sigma^2}{\sigma^2+\lambda^2}.
%\end{equation*}
%\end{p}
%\begin{s}
%Можем да приемем, че $\sigma=1$(ако е с различна стойност от $1$ можем да умножим/разделим $\sigma$ и $\lambda$ до достигане на $\sigma=1$). Нека приемем, че $P[X\geq\lambda]>\dfrac{1}{1+\lambda^2}$. Тъй като математическото очакване за $X$ e $0$,
%\begin{equation*}
%P[X\geq0]E[X|X\geq0]+P[X<0]E[X|X<0]=0
%\end{equation*}
%Знаем също, че
%\begin{equation*}
%P[X\geq0]E[X|X\geq0]>\dfrac{\lambda}{1+\lambda^2}
%\end{equation*}
%и
%\begin{equation*}
%$P[X<0]\leq 1-P[X\geq0]<\dfrac{\lambda^2}{\lambda^2+1}
%\end{equation*}
%От математическото очакване за $X$
%\begin{equation*}
 %P[X<0]E[X|X<0]<-\dfrac{\lambda}{1+\lambda^2}
%\end{equation*}
%\begin{equation*}
%E[X|X<0]<-\dfrac{1}{\lambda}
%\end{equation*}
%Използваме от условието и от свойството на дисперсията, че за всяка случайна величина $X$ $E[X^2]\leq E[X]^2$.
%\begin{equation*}
%E[X^2]\leq P[X\leq\lambda]E[X^2|X\leq\lambda]+P[X<0]E[X^2|X<0]>\dfrac{\lambda^2}{\lambda^2+1}+\dfrac{1}{\lambda^2+1}=1
%\end{equation*}
%Последното се явява в противоречие с факта, че %$E[X^2]=\sigma^2=1$, и така завършваме нашето доказателство.
%\end{s}
\begin{p}[USAMO 2012/2]
Окръжност е разделена на $432$ равни дъги от $432$ точки. Точките са оцветени в $4$ цвята, така че $108$ от тях са оцветени в червено, $108$--в зелено, $108$--в синьо, а останалите $108$ са оцветени в жълто. Докажете, че могат да бъдат избрани по три точки от всеки цвят, така че триъгълниците, определени от едноцветните точки, са еднакви. 
\end{p}
\begin{s}
Ако завъртим червените точки $431$ пъти, те ще се препокрият със зелените точки $108^2$ пъти, което прави средно по $\dfrac{108^2}{431}>27$. Следователно в някакъв момент поне $28$ червени точки препокриват зелените точки и съществуват два еднакви (червен и зелен) $28$-ъгълника.\\
Ако завъртим тези $28$ червени точки $431$ пъти, те ще се препокрият със сините точки $108\cdot28$ пъти, което прави средно по $\dfrac{108\cdot28}{431}>7$. Следователно в някакъв момент поне $8$ червени точки препокриват сините точки и съществуват три еднакви (червен, зелен и син) $8$-ъгълника.\\
Ако завъртим тези $8$ червени точки $431$ пъти, те ще се препокрият с жълтите точки $108\cdot8$ пъти, което прави средно по $\dfrac{108\cdot8}{431}>2$. Следователно в някакъв момент поне $3$ червени точки препокриват жълтите точки и съществуват четири еднакви (червен, зелен, син, жълт) триъгълника.\\
\end{s}
\begin{p}[IMO 2006 C3 Shortlist]
Нека $S$ е множество от точки в равнината, така че никои три от тях да не лежат на една  права. За всеки изпъкнал многоъгълник $P$, означаваме с $a(P)$ броят на върховете на $P$, а с $b(P)$ броят на точките от $S$ извън $P$. Докажете, че за всяко реално число $x$
\begin{equation*}
\sum_{P}x^{a(P)}(1-x)^{b(P)}=1,
\end{equation*}
където сумата е изчислена за всички изпъкнали многоъгълници с върхове във $S$.(Смята се, че отсечка, точка и празното множество са съответно многоъгълници с $2$, $1$ и $0$ върха.)
\end{p}
\begin{s}
Броят на точките в многоъгълника $P$ е $c(P)=|S|-a(P)-b(P)$. Ако положим $y=1-x$ и заместим в лявата част на исканото равенство получаваме
\begin{equation*}
\sum_{P}x^{a(P)}(1-x)^{b(P)}=\sum_{P}x^{a(P)}y^{b(P)}(x+y)^{c(P)}=
\end{equation*}
\begin{equation*}
=\sum_{P}\sum_{i=0}^{c(P)}{{c(P)}\choose{i}}x^{a(P)+i}y^{b(P)+c(P)-i}=
\end{equation*}
\begin{equation*}
=\sum_{P}\sum_{k=a(P)}^{a(P)+c(P)}{{c(P)}\choose{k-a(P)}}x^{k}y^{|S|-k}=
\end{equation*}
\begin{equation*}
={{|S|}\choose{k}}x^{k}y^{|S|-k}=1
\end{equation*}
Коефициентът пред $x^{k}y^{|S|-k}$ е сумата $\sum_{P}{{c(P)}\choose{k-a(P)}}$, която е равна на броя на двойките $(P,H)$ от изпъкнал многоъгълник $P$ и $k$-елементно подмножество от $S$, чиято изпъкнала обвивка е $P$, което пък от своя страна е равно на ${{|S|}\choose{k}}$. Оттук доказателството следва директно. 
\end{s}
\section{Заключение и бъдещо развитие}
В настоящата разработка въведохме един от най-широкоразпространените и полезни комбинаторни методи -- вероятностния метод, и доказахме голямото му приложение. Разгледахме няколко помощни теореми. Бяха дадени примери за основни задачи и бяха обяснени ясно и разбираемо принципите на вероятностния метод. Като извод забелязахме, че наличието на метода в решението на определена задача спомага за простотата му и улеснява решаващия значително.    \\
Като бъдещо развитие авторът си поставя за цел да изследва по-конкретни и/или отворени проблеми и варианти за решаването им чрез вероятностен метод.

\section{Благодарности}

Искам да благодаря на научния си ръководител Димитър Чакъров за това, че ми възложи тази тема и ми предостави нужната литература. Изказвам своите благодарности към него и за конструктивната критика и всички полезни съвети, които ми даде. Благодаря на Никола Стайков за помощта относно текста на разработката.

\begin{thebibliography}{}
\bibitem{2}
Noga Alon, Joel H. Spencer, \textbf {The probabilistic method}, John Wiley and Sons, 2008
\bibitem{2}
Evan Chen, \textbf{Expected uses of probability}, August 2014
\bibitem{2}
Pranav A. Sriram, \textbf{Olympiad Combinatorics}, August 2014
\end{thebibliography}
\end{document}
